\color{unidarkgreen}{\subsection{Логика очувствления}}
\color{black}


\begin{enumerate}

\item
В устройстве предусмотрено два режима очувствления ДЗШ: оперативный и автоматический.
Оперативный режим вводится наличием внешнего сигнала \textit{«Оперативный ввод очувствления»} и применяется для обеспечения требуемой чувствительности ДЗШ при выводе в ремонт мощных присоединений.
Автоматический режим вводится наличием внешнего сигнала \textit{«Нормальный режим очувствления»} или от программного ключа \textbf{«Очувст.»} и предназначен для повышения надежности отключения выключателей присоединений при работе ДЗШ, в том числе при АПВ шин и при опробовании.

\item
Режим очувствления, вводимый срабатыванием ДЗШ, сохраняется на всё время отключения СШ.
При наличии сигнала \textit{«Фикс. откл 1С»} (\textit{«Фикс. откл. 2С»}) режим очувствления удерживается до момента возможного повторного включения СШ при АПВ.
Сброс очувствления до повторного включения осуществляется при наличии сигнала \textit{«Работа АПВ 1С»} (\textit{«Работа АПВ 2С»}).

\item
Переход в режим очувствления при опробовании осуществляется при отсутствии напряжения на СШ.
Очувствление в данном случае сохраняется на всё время опробования.

\item
При АПВ шин режим очувствления вводится на время включения выключателя и последующего отключения в случае неуспешного АПВ.
Для исключения излишнего очувствления при постепенном включении всех выключателей предусмотрена возможность вывода режима очувствления при успешном АПВ по истечению выдержки времени, задаваемой уставкой \textbf{«Точ. АПВ 1С»} (\textbf{«Точ. АПВ 2С»}).
Вывод режима осуществляется программным ключом \textbf{«Вывод очув. при успеш. АПВ 1С»} (\textbf{«Вывод очув. при успеш. АПВ 2С»}).

\item 
Уставки логики очувствления приведены в таблице \ref{app_set:loch}. Алгоритм работы логики очувствления приведен на рисунке \ref{pic:Log_ochuv}.

\medskip
\begin{figure}[H]
\centering
\includegraphics[width=0.9\textwidth,height=0.9\textheight,keepaspectratio]{\fbpath//04.01 ДЗШ/661.1771 Логика очувствления/img/Log_ochuv.pdf}
\captionof{figure}{Функциональная схема логики очувствления}\label{pic:Log_ochuv}
\end{figure}



\end{enumerate}
